\chapter{Introduction and Motivation}
\label{ch:introduction}

\section{Motivation}
\label{sec:motivation}

A few seconds of a film score are often enough to tell what kind of film it belongs to.
A sustained dissonant string cluster over a low pulse points to a horror film, and a brass
fanfare over driving percussion to an action film. A solo piano over soft strings suggests
a drama or a romance. Listeners make this judgement without seeing a single frame, and
they make it with some confidence. The starting point of this thesis is that the judgement
is not really about genre at all. A listener first perceives an emotion (threat,
excitement, tenderness, melancholy) and infers the genre from it.

This inference works because film composers deliberately write music that signals the
emotional register a scene and a film are meant to have. \citet{gorbman1987unheard}
describes film music as a system of learned codes. Under this view, minor keys and
dissonant harmony conventionally signal threat or suspense, while other harmonic,
rhythmic and instrumental choices are conventionally tied to romance, comedy or heroism.
\citet{juslin2013everyday} places such codes within a general account of how music comes
to carry emotional meaning, in which learned associations let music evoke culturally
shared emotional connotations. In film scoring, these connotations are closely tied to
genre convention. Section~\ref{sec:fund_motivation_theory} discusses both accounts in
detail. For this thesis, it follows that the link between a film's soundtrack and its
genre is expected to run through emotion, because composers place it there.

Automatic genre prediction from film audio is an established task with practical value.
Genre metadata organises film archives and streaming catalogues, drives recommendation,
and supports the search for music that fits a scene. Existing work, however, treats the
task as a pure pattern-recognition problem. \citet{austin2010characterization} classify
film genre from low-level features of the score, \citet{ma2021computational} from learned
audio embeddings of the music cues used in a film, and \citet{mangolin2020multimodal} and
\citet{behrouzi2023multimodal} from trailer audio combined with visual and textual
channels. These studies are reviewed in Chapter~\ref{ch:relatedwork}. In all of them,
genre is predicted \emph{directly} from an audio representation. Such models may be
accurate, but they cannot explain why a soundtrack sounds like Horror. Their internal
representation is a vector of several hundred numbers without a fixed meaning, and
nothing in the model corresponds to the emotional mechanism described by composers and
film-music scholars.

This leaves a gap between what is known about film music and how it is modelled. If
emotion is the mechanism that connects a soundtrack to its genre, then modelling the
emotion explicitly should cost little accuracy compared with modelling around it, and
the result should be much easier to interpret. A model that predicts Horror because the
music is fearful and tense can be checked against film-music practice, questioned and
corrected by a human. This is not possible if the reason is that dimension 87 of an
embedding is large. There is also a second, less obvious reason to expect an emotional
representation to help. An emotion rating means the same thing in any dataset, since a
fearful cue is fearful whichever film collection it comes from. The dimensions of a
learned audio embedding instead carry the statistics of the data they were computed on.
A model built on emotions might therefore survive a change of dataset better than a
model built on raw embeddings.

Neither expectation is guaranteed. Emotion is a strong compression of the audio, and a
compression can only help if what it keeps matters more than what it discards. The
emotions must also be predicted from audio before they can be used, and emotion
prediction from audio is far from perfect. Whether an explicit emotional intermediate
helps, costs nothing or simply loses information is therefore an empirical question.
This thesis addresses it by asking: \emph{can the emotional features of a film
soundtrack predict the genre of the film, and does modelling emotion as an interpretable
intermediate step do better than classifying genre directly from audio?}

\section{Aim of this thesis}
\label{sec:aim}

The aim is to test whether soundtrack emotion is a usable and interpretable route from
film audio to film genre. The work is organised around four research questions, and
Chapter~\ref{ch:evaluation} states which of its sections answers each.

\begin{description}
  \item[RQ1] Can a soundtrack's emotional profile predict its film's genre?
  \item[RQ2a] Which emotions relate to which genres?
  \item[RQ2b] Does routing audio through an emotion intermediate beat classifying genre
    directly from audio?
  \item[RQ3] Which emotions are the strongest discriminators for each genre?
\end{description}

RQ1 asks whether the signal exists at all, first with emotions rated by human listeners
and then with emotions predicted from audio. RQ2a and RQ3 concern the form of the
signal. RQ2a maps the relationship between every emotion and every genre, and RQ3
identifies for each genre the emotions that distinguish it from the rest. Together they
test whether the intermediate can be read in the way the motivation above argues. RQ2b
is the central comparative question of the thesis. It is asked within a single dataset
and across datasets, where a model trained on one collection of films is applied to
another without seeing any of its labels.

\mytodo{To be written by the author: the list of contributions. A draft is kept outside the project.}

\section{Structure of this thesis}
\label{sec:structure}

Chapter~\ref{ch:fundamentals} explains, for readers new to the field, the models of
emotion, the audio representations, the two learning methods and the evaluation
apparatus used in this thesis. Chapter~\ref{ch:relatedwork} reviews the surveys of music
emotion recognition, the evidence for an emotion--genre relationship in music, and the
four film-genre studies this work is compared with. Chapter~\ref{ch:methods} describes
the two datasets, the genre labels, the two-stage pipeline, the baselines and controls,
and the evaluation protocol. Chapter~\ref{ch:evaluation} reports the experiments,
grouped by research question. Chapter~\ref{ch:discussion} interprets the results and
discusses the limitations of the work, and Chapter~\ref{ch:conclusions} answers the
research questions and names the next steps.
